The strongest general theorem currently available here is the regularization theorem of Tehrani–McLean–Zinger: after a deformation through cohomologous symplectic forms preserving the SC divisor, one gets a regularized neighborhood model of the divisor, and their later work formalizes the associated “normal-direction integrable” almost-complex structures. ([math.stonybrook.edu](https://www.math.stonybrook.edu/~azinger/research/SympDivConf.pdf))

Using that input, one can give a complete proof of the **full local holomorphic-coordinate statement in real dimension four**. The proof below is fully rigorous in that case. The only place where the argument is dimension-specific is the automatic integrability of a \((0,1)\)-connection on a line bundle over a complex curve.

```latex
\documentclass[11pt]{article}

\usepackage[T1]{fontenc}
\usepackage{amsmath,amsthm,amssymb,mathtools}
\usepackage{enumitem}
\usepackage{hyperref}
\usepackage[margin=1in]{geometry}

\newtheorem{theorem}{Theorem}
\newtheorem{proposition}[theorem]{Proposition}
\newtheorem{lemma}[theorem]{Lemma}
\newtheorem{corollary}[theorem]{Corollary}
\theoremstyle{remark}
\newtheorem{remark}[theorem]{Remark}

\newcommand{\RR}{\mathbb{R}}
\newcommand{\CC}{\mathbb{C}}

\title{Local Holomorphic Equations for Symplectic Simple Normal Crossings Divisors in Dimension Four}
\author{}
\date{}

\begin{document}
\maketitle

\begin{abstract}
We prove the following four-dimensional case of the question posed in the prompt.

\smallskip

\noindent
\textbf{Main theorem.}
\emph{Let $(M^4,\omega)$ be a closed symplectic $4$--manifold and
$D\subset M$ a symplectic simple normal crossings divisor.
After replacing $\omega$ by a cohomologous symplectic form through a path of
symplectic forms for which $D$ remains a symplectic simple normal crossings divisor,
there exists an $\omega$--tame almost complex structure $J$ such that for every
$p\in M$ there are smooth complex coordinates $(z_1,z_2)$ centered at $p$
with
\[
D=\{z_1\cdots z_r=0\}
\qquad\text{for some }r\in\{0,1,2\},
\]
and $z_1,\dots,z_r$ are $J$--holomorphic.}

\smallskip

The proof uses the regularization theorem of Tehrani--McLean--Zinger for
symplectic simple normal crossings divisors and then upgrades the local models
to genuine holomorphic normal coordinates by exploiting the special feature that
every irreducible component of $D$ is a symplectic surface.
\end{abstract}

\section{Statement and preliminary input}

Write
\[
D=\bigcup_{i\in S} D_i
\]
as the union of its smooth irreducible components.
Since $\dim_{\RR}M=4$ and the components have codimension $2$, transversality implies:

\begin{lemma}\label{lem:doublepoints}
Every intersection $D_i\cap D_j$ is a finite set, and there are no triple intersections.
\end{lemma}

\begin{proof}
A transverse intersection of two codimension-$2$ submanifolds in a
$4$--manifold is $0$--dimensional, hence discrete; compactness gives finiteness.
A transverse intersection of three codimension-$2$ submanifolds in a $4$--manifold
would have expected dimension $4-2-2-2=-2$, which is impossible.
\end{proof}

The only external input in the proof is the following consequence of the
Tehrani--McLean--Zinger regularization theorem.

\begin{proposition}[Tehrani--McLean--Zinger]\label{prop:TMZ}
After deforming $\omega$ through cohomologous symplectic forms for which $D$
remains a symplectic simple normal crossings divisor, we may assume that $D$
admits an $\omega$--regularization in the sense of
\cite[Definition~2.12]{TMZ}.
In the present four-dimensional situation this means:

\begin{enumerate}[label=\textup{(\roman*)}]
\item
for each component $D_i$ there is a complex line bundle structure on
its symplectic normal bundle
\[
N_i:=N_M D_i \to D_i,
\]
a radius function $\rho_i$, a connection $1$--form $\alpha_i$,
and a tubular neighborhood embedding
\[
\Psi_i\colon \mathcal N_i \to M
\]
from a neighborhood $\mathcal N_i$ of the zero section in $N_i$, such that
$\Psi_i|_{D_i}=\mathrm{id}_{D_i}$ and
\[
\Psi_i^*\omega
=
\pi_i^*(\omega|_{D_i})+\frac12\, d(\rho_i\alpha_i);
\]

\item
for each node $p\in D_i\cap D_j$ there is a neighborhood embedding
\[
\Psi_p\colon B_p\to M
\]
from a neighborhood $B_p\subset (N_i)_p\oplus (N_j)_p\cong\CC^2$ of the origin
such that
\[
\Psi_p(0)=p,\qquad
\Psi_p^{-1}(D)=B_p\cap\bigl(\{z_i=0\}\cup\{z_j=0\}\bigr),
\]
for the standard complex coordinates $(z_i,z_j)$ on $(N_i)_p\oplus (N_j)_p$,
and such that these nodal charts are compatible with the component charts
$\Psi_i,\Psi_j$ in the sense encoded in the definition of regularization.
\end{enumerate}
\end{proposition}

\begin{proof}
This is exactly the specialization to the case of a four-dimensional SC divisor
of \cite[Theorem~2.13]{TMZ}; see also the summary in \cite[\S3]{TMZCRAS}.
\end{proof}

\begin{remark}
The proposition is the all-dimensional part of the argument.  The remainder of the proof
is four-dimensional and uses that each $D_i$ is a surface.
\end{remark}

\section{A holomorphic line-bundle lemma over a Riemann surface}

We isolate the elementary complex-analytic fact that replaces the missing
higher-dimensional integrability.

\begin{lemma}\label{lem:holomorphic_line_bundle}
Let $\Sigma$ be a smooth oriented surface endowed with an integrable complex
structure $j$, and let $L\to\Sigma$ be a smooth complex line bundle.
Fix finitely many pairwise disjoint coordinate discs
\[
U_a\subset \Sigma
\]
and on each $U_a$ fix a smooth complex trivialization
\[
\Phi_a\colon L|_{U_a}\xrightarrow{\;\sim\;} U_a\times \CC.
\]
Then there exists a holomorphic structure on $L$ for which, after possibly shrinking
the discs $U_a$, each chosen trivialization $\Phi_a$ becomes holomorphic.
\end{lemma}

\begin{proof}
Choose a smooth complex connection $\nabla^0$ on $L$.
Because the discs are pairwise disjoint, we may choose smaller discs
$U'_a\Subset U_a$ and a partition of unity to modify $\nabla^0$ so as to obtain a
new complex connection $\nabla$ with the property that, in the fixed trivialization
over $U'_a$, the connection form of $\nabla$ vanishes identically.
(Indeed, in the trivialization over $U_a$ we may write
$\nabla^0=d+A_a$, where $A_a$ is a complex-valued $1$--form; multiplying $A_a$
by a bump function equal to $1$ on $U'_a$ and subtracting the resulting sum from
$\nabla^0$ produces $\nabla$.)

We now claim that the $(0,1)$--part
\[
\bar\partial_\nabla := \nabla^{0,1}
\]
defines a holomorphic structure on $L$.
Since $\Sigma$ has complex dimension $1$, there are no $(0,2)$--forms, hence
$(\bar\partial_\nabla)^2=0$ automatically.

It remains to show local existence of nonvanishing $\bar\partial_\nabla$--holomorphic
frames.
Let $V\subset\Sigma$ be any complex coordinate disc with coordinate $\zeta$ and
let $e$ be a smooth nowhere-vanishing local frame of $L$ over $V$.
Then
\[
\bar\partial_\nabla e = a\, e\, d\bar\zeta
\]
for some smooth complex-valued function $a$ on $V$.
A section $s=f e$ is $\bar\partial_\nabla$--holomorphic iff
\[
\frac{\partial f}{\partial\bar\zeta}+a f=0.
\]
Choose a smooth function $u$ on a smaller disc $V'\Subset V$ solving
\[
\frac{\partial u}{\partial\bar\zeta}=-a.
\]
(Such a solution exists by the classical local solvability of the scalar
$\bar\partial$--equation on planar domains, e.g.\ via the Cauchy--Green formula.)
Then $f=e^u$ satisfies
\[
\frac{\partial f}{\partial\bar\zeta}+af
=
\left(\frac{\partial u}{\partial\bar\zeta}+a\right)e^u=0,
\]
so $s=e^u e$ is a nonvanishing local $\bar\partial_\nabla$--holomorphic frame.

Thus $\bar\partial_\nabla$ defines a holomorphic line-bundle structure on $L$.
By construction, on each $U'_a$ the connection form of $\nabla$ vanishes in the
prescribed trivialization, so the standard frame is $\bar\partial_\nabla$--holomorphic
there.  Hence the chosen trivialization is holomorphic on $U'_a$.
\end{proof}

\section{Proof of the main theorem}

\begin{theorem}\label{thm:main}
Let $(M^4,\omega)$ be a closed symplectic $4$--manifold and
$D\subset M$ a symplectic simple normal crossings divisor.
After replacing $\omega$ by a cohomologous symplectic form through a path of
symplectic forms for which $D$ remains a symplectic simple normal crossings divisor,
there exists an $\omega$--tame almost complex structure $J$ such that for every
$p\in M$ there are smooth complex coordinates $(z_1,z_2)$ centered at $p$
with
\[
D=\{z_1\cdots z_r=0\}
\qquad\text{for some }r\in\{0,1,2\},
\]
and $z_1,\dots,z_r$ are $J$--holomorphic.
\end{theorem}

\begin{proof}
By Proposition~\ref{prop:TMZ}, after the allowed deformation of $\omega$ we may assume
that $D$ admits an $\omega$--regularization.
We keep the same notation $\omega$ for the deformed symplectic form.

\medskip
\noindent\textbf{Step 1: choose complex structures on the components.}
For each component $D_i$, let
\[
P_i:= D_i\cap \bigcup_{j\neq i} D_j
\]
be the finite set of nodes on $D_i$.
For each $p\in P_i$, choose a nodal regularization chart
\[
\Psi_p\colon B_p\to M
\]
as in Proposition~\ref{prop:TMZ}, with complex coordinates $(z_i,z_j)$ on
$B_p\subset\CC^2$ such that
\[
D_i\cap \mathrm{Im}(\Psi_p)=\Psi_p(B_p\cap\{z_i=0\}),\qquad
D_j\cap \mathrm{Im}(\Psi_p)=\Psi_p(B_p\cap\{z_j=0\}).
\]
After shrinking the $B_p$, we may assume that the images $\Psi_p(B_p)$ are pairwise disjoint.
The coordinate $z_j$ restricts to a complex coordinate on the branch $D_i$
near the point $p$.

We now choose on each surface $D_i$ a global almost complex structure $j_i$
compatible with $\omega|_{D_i}$ and equal, on a neighborhood of each node
$p\in P_i$, to the standard complex structure determined by the local coordinate
$z_j$ above.

To do this, choose any background Riemannian metric $g_i^0$ on $D_i$, and on each
nodal disc choose the Euclidean metric in the coordinate $z_j$.
Using a partition of unity, glue these to a global metric $g_i$ on $D_i$ that is
Euclidean in the nodal coordinates near each node.
By the standard linear algebra of symplectic $2$--planes, the pair
$(\omega|_{D_i},g_i)$ determines a unique $\omega|_{D_i}$--compatible almost
complex structure $j_i$.
Because $D_i$ is a surface, every almost complex structure on it is integrable.
Hence $(D_i,j_i)$ is a Riemann surface, and near each node the coordinate $z_j$
is $j_i$--holomorphic.

\medskip
\noindent\textbf{Step 2: put holomorphic structures on the normal line bundles.}
Fix $i\in S$.
Let $N_i=N_M D_i\to D_i$ be the complex line bundle coming from the regularization.
For each node $p\in P_i$, the compatibility of the regularization identifies
$N_i$ over a small $j_i$--holomorphic disc $U_{i,p}\subset D_i$ with the trivial bundle
\[
U_{i,p}\times \CC_{z_i},
\]
the base coordinate being the branch coordinate $z_j$ and the fiber coordinate
being the $z_i$--coordinate of the nodal chart.
Shrinking the $U_{i,p}$ if necessary, we may assume that they are pairwise disjoint.

Apply Lemma~\ref{lem:holomorphic_line_bundle} to the complex line bundle
$N_i\to (D_i,j_i)$ with the prescribed trivializations above.
We obtain a holomorphic structure on $N_i$ for which, on smaller discs
$U'_{i,p}\Subset U_{i,p}$, the trivializations
\[
N_i|_{U'_{i,p}} \cong U'_{i,p}\times \CC_{z_i}
\]
are holomorphic.

Let $J_i^{\mathrm{hol}}$ denote the resulting integrable complex structure on
the total space of $N_i$ near the zero section.
By construction:

\begin{itemize}
\item
the zero section $D_i\subset N_i$ is a complex curve;
\item
near each node $p\in P_i$, in the coordinates $(z_j,z_i)$ the complex structure
$J_i^{\mathrm{hol}}$ is the standard product complex structure on
$U'_{i,p}\times \CC$.
\end{itemize}

\medskip
\noindent\textbf{Step 3: tame the regularized symplectic form near each component.}
Consider the symplectic form $\Psi_i^*\omega$ on a neighborhood of the zero section in $N_i$.
At points of the zero section, the defining property of the regularization gives the splitting
\[
T(N_i)|_{D_i}=TD_i\oplus N_i
\]
and
\[
\Psi_i^*\omega|_{D_i}
=
\omega|_{D_i}\oplus \omega_{\mathrm{fib}},
\]
where $\omega_{\mathrm{fib}}$ is the fiberwise symplectic form on the complex line
bundle $N_i$.
On the zero section, $J_i^{\mathrm{hol}}$ acts by $j_i$ on $TD_i$ and by multiplication
by $i$ on the fibers; therefore
\[
\Psi_i^*\omega(v, J_i^{\mathrm{hol}}v)>0
\qquad
\forall\, v\neq 0
\]
along the zero section.

Since positivity is an open condition and the zero section is compact, after shrinking
the neighborhood in $N_i$ we may assume that $\Psi_i^*\omega$ tames
$J_i^{\mathrm{hol}}$ on that whole neighborhood.
Let
\[
\mathcal U_i\subset N_i
\]
be such a shrunk neighborhood.

\medskip
\noindent\textbf{Step 4: glue the component models with the nodal models.}
For each node $p\in D_i\cap D_j$, shrink $B_p$ if necessary so that:
\begin{enumerate}[label=\textup{(\alph*)}]
\item
$\Psi_p(B_p)$ meets no nodal chart except itself;
\item
$\Psi_p(B_p)$ is contained in $\Psi_i(\mathcal U_i)\cup \Psi_j(\mathcal U_j)$;
\item
in the overlaps with $\Psi_i(\mathcal U_i)$ and $\Psi_j(\mathcal U_j)$,
the complex structures coming from $J_i^{\mathrm{hol}}$, $J_j^{\mathrm{hol}}$,
and the standard complex structure on $B_p\subset\CC^2$ agree.
\end{enumerate}
Condition \textup{(c)} holds because, by construction in Step~2, the component
complex structures are literally the standard product structure in the coordinates
$(z_j,z_i)$ and $(z_i,z_j)$ near the node.

Thus the pushforwards of the $J_i^{\mathrm{hol}}$ and of the standard complex structures
on the nodal charts glue to a well-defined smooth almost complex structure
\[
J_{\mathrm{loc}}
\]
on some open neighborhood $U$ of $D$ in $M$.
By construction, $J_{\mathrm{loc}}$ is $\omega$--tame on $U$.

\medskip
\noindent\textbf{Step 5: extend globally.}
Let $\mathcal J_\tau(TM,\omega)\to M$ be the bundle whose fiber at $x\in M$
is the contractible space of $\omega_x$--tame complex structures on $T_xM$.
The local almost complex structure $J_{\mathrm{loc}}$ is a smooth section of this
bundle over the open neighborhood $U$ of the closed set $D$.
Since the fibers are contractible, standard obstruction theory for smooth bundles
extends this section to a global smooth section
\[
J\in \Gamma(\mathcal J_\tau(TM,\omega)).
\]
After shrinking $U$ once more if necessary, we may assume that $J=J_{\mathrm{loc}}$
on $U$.
In particular, $J$ is $\omega$--tame on all of $M$.

\medskip
\noindent\textbf{Step 6: verify the coordinate description.}
Let $p\in M$.

\smallskip
\emph{Case 1: $p\notin D$.}
Take any smooth complex coordinates $(z_1,z_2)$ near $p$; then $r=0$.

\smallskip
\emph{Case 2: $p\in D_i\setminus \bigcup_{j\neq i}D_j$.}
Since $J=J_{\mathrm{loc}}$ near $p$, the point lies in the image of one component
chart $\Psi_i(\mathcal U_i)$ and away from the nodal region.
Because $N_i$ is a holomorphic line bundle over the Riemann surface $(D_i,j_i)$,
there exist local holomorphic coordinates $(\zeta,w)$ on the total space of $N_i$
with zero section given by $w=0$.
Pushing forward by $\Psi_i$, we obtain smooth complex coordinates
\[
(z_1,z_2)=(w,\zeta)
\]
near $p$ in which $z_1$ is $J$--holomorphic and
\[
D=\{z_1=0\}.
\]
Thus $r=1$.

\smallskip
\emph{Case 3: $p\in D_i\cap D_j$.}
By Lemma~\ref{lem:doublepoints}, no third component passes through $p$.
Since $J=J_{\mathrm{loc}}$ near $p$, the nodal chart $\Psi_p$ identifies a
neighborhood of $p$ with a neighborhood of the origin in $\CC^2$ endowed with
the standard complex structure and coordinates $(z_i,z_j)$.
In these coordinates,
\[
D=\{z_i z_j=0\},
\]
and both $z_i$ and $z_j$ are $J$--holomorphic.
Thus $r=2$.

\smallskip
These three cases prove the theorem.
\end{proof}

\begin{remark}
The proof is genuinely four-dimensional.
The only place where this enters is Step~2:
over a complex curve, a $(0,1)$--connection on a line bundle is automatically
integrable because there are no $(0,2)$--forms.
In higher dimensions this automatic integrability fails, so the argument above does
not extend verbatim.
\end{remark}

\begin{thebibliography}{99}

\bibitem{MS}
D. McDuff and D. Salamon,
\emph{J-holomorphic Curves and Symplectic Topology},
2nd ed.,
American Mathematical Society Colloquium Publications, Vol.~52,
Amer. Math. Soc., Providence, RI, 2012.

\bibitem{TMZ}
M. F. Tehrani, M. McLean, and A. Zinger,
\emph{Normal crossings singularities for symplectic topology},
Adv. Math. \textbf{339} (2018), 672--748.

\bibitem{TMZCRAS}
M. F. Tehrani, M. McLean, and A. Zinger,
\emph{Singularities and semistable degenerations for symplectic topology},
C. R. Acad. Sci. Paris Ser. I \textbf{356} (2018), 420--432.

\bibitem{TMZStr}
M. Farajzadeh-Tehrani, M. McLean, and A. Zinger,
\emph{Normal Crossings Singularities for Symplectic Topology: Structures},
Acta Math. Sin. (Engl. Ser.) \textbf{40} (2024), no.~1, 107--160.

\end{thebibliography}

\end{document}
```

If you want, I can also produce a second version that is optimized for publication style (leaner introduction, more formal bibliography, fewer explanatory remarks).